% MRD Stage 0--2: the core answers to the 2026-10-08 questions (hard limit 12 pages, target 10).
% Build: bash build.sh (the default target: figures, then report_core.pdf). See README.md.
%
% This is the short read. The long companion is report.tex (built as report_full.pdf), whose section numbers
% the "Full report" pointers on every page use. Nothing is computed for this file: every number is copied from
% the full report's sections/*.tex or from the devlog entry named in a page's [Source: ...] pointer.
% One file per page in sections_core/; drop a page by commenting out its \input line below.
% Figures are reused from figures/ (written by the figs_NN_*.py modules). Two are cropped here with
% \includegraphics trim (fig02a, fig02b_d26a); if their layout changes, check the crops (README.md).

\documentclass[10pt]{article}
% Shared preamble for the 2026-10-08 questions report. Section files use only these macros.
\usepackage[T1]{fontenc}
\usepackage[utf8]{inputenc}
\usepackage{lmodern}
\usepackage{microtype}
\usepackage[letterpaper,margin=1in]{geometry}
\usepackage{graphicx}
\usepackage{booktabs}
\usepackage{tabularx}
\usepackage{array}
\usepackage{longtable}
\usepackage{multirow}
\usepackage{amsmath}
\usepackage{amssymb}
\usepackage{enumitem}
\usepackage{caption}
\usepackage{float}
\usepackage{parskip}
\usepackage[dvipsnames]{xcolor}
\usepackage{tcolorbox}
\usepackage{fancyhdr}
\usepackage{titlesec}
\usepackage[hidelinks]{hyperref}

\definecolor{inkmuted}{HTML}{52514E}
\definecolor{boxtake}{HTML}{EAF2FC}
\definecolor{boxtakeedge}{HTML}{2A78D6}
\definecolor{boxcaveat}{HTML}{F4F3EF}
\definecolor{boxcaveatedge}{HTML}{898781}
\definecolor{boxteam}{HTML}{FDF0EA}
\definecolor{boxteamedge}{HTML}{EB6834}

\setlist{itemsep=2pt,topsep=3pt,parsep=0pt}
\captionsetup{font=small,labelfont=bf,margin=0pt,skip=4pt}
\setlength{\parskip}{5pt}

\pagestyle{fancy}
\fancyhf{}
\fancyhead[L]{\small\color{inkmuted}MRD Stage 0--2: answers to the 2026-10-08 questions}
\fancyhead[R]{\small\color{inkmuted}\thepage}
\renewcommand{\headrulewidth}{0.3pt}

\titleformat{\section}{\Large\bfseries}{\thesection}{0.6em}{}
\titleformat{\subsection}{\large\bfseries}{\thesubsection}{0.6em}{}

% The short answer to the section's question (one per section, at the top).
\newtcolorbox{answerbox}{colback=boxtake,colframe=boxtakeedge,boxrule=0.6pt,arc=2pt,left=6pt,right=6pt,top=4pt,bottom=4pt,fonttitle=\bfseries\small,title=Short answer}
% What this result does NOT show (claim boundary).
\newtcolorbox{caveatbox}{colback=boxcaveat,colframe=boxcaveatedge,boxrule=0.6pt,arc=2pt,left=6pt,right=6pt,top=4pt,bottom=4pt,fonttitle=\bfseries\small,coltitle=black,title=What this does not show}
% A sentence the user can say to the team, and the likely follow-up question.
\newtcolorbox{teambox}{colback=boxteam,colframe=boxteamedge,boxrule=0.6pt,arc=2pt,left=6pt,right=6pt,top=4pt,bottom=4pt,fonttitle=\bfseries\small,title=For the team}
% A concept explainer (intuition first, then the formula). Drawn as a left rule with a plain title, so it
% cannot be mistaken for the framed gray "What this does not show" box.
\newtcolorbox{explainbox}[1]{colback=white,colframe=inkmuted!55,colbacktitle=white,coltitle=black,boxrule=0pt,leftrule=2.5pt,titlerule=0pt,arc=0pt,left=8pt,right=6pt,top=3pt,bottom=4pt,toptitle=2pt,bottomtitle=0pt,before upper={\setlength{\parskip}{3pt}},fonttitle=\bfseries\small,title={#1}}

% Source pointer: devlog id / analysis id, set small and muted.
\newcommand{\src}[1]{{\small\color{inkmuted}[Source: #1]}}
\newcommand{\pp}{\,pp}
\newcommand{\lift}{\mathrm{lift}}


% Tighter page than the full report: the core answers are meant to fit about one page each.
\geometry{letterpaper,left=0.7in,right=0.7in,top=0.75in,bottom=0.65in,headsep=10pt,footskip=18pt}
\setlength{\parskip}{3pt}
\setlist{itemsep=1.5pt,topsep=2pt,parsep=0pt}
\captionsetup{font=footnotesize,labelfont=bf,margin=0pt,skip=3pt}
\fancyhead[L]{\small\color{inkmuted}MRD Stage 0--2: core answers to the 2026-10-08 questions (v3.3, 2026-10-11)}
\titleformat{\section}{\large\bfseries}{}{0em}{}
\titlespacing*{\section}{0pt}{0pt}{4pt}
\renewcommand{\thefigure}{\arabic{figure}}

\hypersetup{
  pdftitle={MRD Stage 0--2: core answers to the 2026-10-08 questions},
  pdfsubject={Short read for the project lead and collaborators; descriptive results, seven supported patients},
}

% ---- page macros (core only)
% One question per page: \coreq{title}{label}{the question in the asker's words}
\newcommand{\coreq}[3]{%
  \clearpage
  \phantomsection\section*{#1}\label{#2}%
  {\small\color{inkmuted}\emph{Asked:} #3\par}\vspace{2pt}}
% Half-page item that does not force a new page.
\newcommand{\corehalf}[3]{%
  \phantomsection\section*{#1}\label{#2}%
  {\small\color{inkmuted}\emph{Asked:} #3\par}\vspace{2pt}}
% Evidence list heading.
\newcommand{\evidence}{\noindent\textbf{Evidence}\par\nobreak\vspace{-1pt}}
% One-line claim boundary.
\newcommand{\doesnot}[1]{%
  \begin{tcolorbox}[colback=boxcaveat,colframe=boxcaveatedge,boxrule=0.4pt,arc=1pt,left=5pt,right=5pt,top=2pt,bottom=2pt,before skip=4pt,after skip=3pt]
  \small\textbf{Does not show:} #1
  \end{tcolorbox}}
% Pointer to the long companion.
\newcommand{\fullref}[1]{{\small\color{inkmuted}\textbf{Detail:} full report (\texttt{report\_full.pdf}), #1\par}}
% Compact short-answer box (same colours as the full report's).
\newtcolorbox{coreanswer}{colback=boxtake,colframe=boxtakeedge,boxrule=0.6pt,arc=2pt,left=6pt,right=6pt,top=3pt,bottom=3pt,before skip=3pt,after skip=5pt,fonttitle=\bfseries\small,title=Short answer}

\begin{document}

% Core page 1: title block and the answers-at-a-glance table. Page numbers resolve through \pageref;
% the last column is the section of the full report (report_full.pdf) that holds the detail.
\thispagestyle{plain}
\begin{flushleft}
{\LARGE\bfseries MRD Stage 0--2: the core answers\par}
\vspace{3pt}
{\large\color{inkmuted}The 2026-10-08 questions, one page each \textperiodcentered{} version 3.3, 11 October 2026\par}
\end{flushleft}
\vspace{-2pt}
{\small
This is the short read. Each page answers one question with a short answer, 3--5 pieces of evidence with their numbers, one figure, a line on what the result does \emph{not} show, and a pointer to the section of the long companion (\texttt{report\_full.pdf}) that has the detail. Terms are defined in a few words where they first appear, and in a box on page~\pageref{core:cohort}.

\emph{MX} is Minhang Xu, the study lead and repository owner: ``MX chose'' marks a decision of the study's, and the last page lists the decisions still open for MX. Every result here is \emph{descriptive}: an effect size measured on these data (seven supported patients; a Stage~1 basis mostly fit on all cells, held-out patient included), with no calibrated test, no patient-level interval and no outside cohort behind it. The study's claim list is THEME \S6 (\texttt{scripts/\allowbreak mrd\_stage0\_2/\allowbreak research/\allowbreak THEME.md}).\par}

\vspace{4pt}
\noindent{\large\bfseries Answers at a glance}\par\vspace{2pt}
{\footnotesize
\setlength{\tabcolsep}{3.5pt}
\renewcommand{\arraystretch}{1.18}
\noindent\begin{tabularx}{\linewidth}{@{}>{\raggedright\arraybackslash}p{3.35cm} >{\raggedright\arraybackslash}X >{\raggedright\arraybackslash}p{3.0cm} >{\centering\arraybackslash}p{0.72cm} >{\centering\arraybackslash}p{0.6cm}@{}}
\toprule
\textbf{Question} & \textbf{Answer} & \textbf{How sure} & \textbf{Page} & \textbf{Full} \\
\midrule
What is the cohort, and how do we score a classifier? &
13 patients, 60,118 cells, 1,569 malignant; 7 supported patients. Lift = held-out AUPRC minus the patient's malignant prevalence $\pi$, in percentage points (pp). &
Counts and definitions &
\pageref{core:cohort} & 1 \\
Do P02, P03, P04 disproportionately shape the factors? (IFN-$\alpha$) &
No more than other patients: without them the IFN-$\alpha$ axis barely moves (matched $|r| \geq 0.988$). Their imbalance is mainly in Stage~2: 68.3\% of the malignant labels. &
Descriptive; leverage on 1 panel, score variance on 41; PCA/FA only &
\pageref{core:shape} & 2 \\
What did we do with allograft rejection and apoptosis? &
Ran the same leverage test (D26). Leaving P02+P03+P04 out moves the axis beyond matched random drops (mostly P03), not beyond other patients on the primary axis; P01 moves it most. &
Descriptive; 2 panels, PCA/FA; rule fixed before the run &
\pageref{core:d26} & 2b \\
Should we aggregate patients? &
Yes, on the current bases: 45 of 47 representations pass all four overlap gates, 2 are location shifts, none looks like a patient-specific subtype. &
Descriptive; transductive basis; gates are judgment calls &
\pageref{core:aggregate} & 3 \\
IFN-$\alpha$: why does PCA40 reach four patients (lift $>$ 5\pp) but FA40 only P03 and P04? &
An L1 entry threshold, not missing FA structure: FA7 misses the entry bar by 1.4--3.3\%; with $C$ tuned per fold FA40 reaches all 7 patients. &
Descriptive; 1 panel; the rule reproduces 2,772/2,772 audited folds &
\pageref{core:pcafa} & 4 \\
What is PC8? &
One axis, named by its poles: malignant cells lean toward TXNIP, IFITM2, TENT5A, normal cells toward IFI27, PSME2, PSMA3. A refit can spread it over several factors. &
Descriptive; 1 panel &
\pageref{core:pc8} & 5 \\
Have we integrated the protein data? Does it track PC8? &
Not into the factors (RNA only); we use it to read them. PC8's strongest correlates (CD34; CD48, CD116, CD33, CD146) hold within patients but not inside cell types; CD99 is weaker within patients yet keeps its size inside them (a lead). &
Descriptive; 1 panel; protein CLR-normalized, no dsb &
\pageref{core:protein} & 5 \\
Exclude the held-out patient from Stage~1 (Option~2)? &
Done on IFN-$\alpha$: lift stays positive in 7/7 patients; the median per-patient cost is 0.05--1.29\pp, so transductive numbers are slightly optimistic. &
Descriptive; IFN-$\alpha$ only, PCA/FA &
\pageref{core:inductive} & 6 \\
AUPRC against the number of cells or patients? &
Drawn, with $C$ retuned: one supported training patient already gives positive lift; going from two to six adds 0.30--2.34\pp, CIs mostly include 0. &
Descriptive; frozen basis; nothing extrapolated &
\pageref{core:curves} & 7 \\
Q1--Q3: shared or patient-specific? how strong? what does DR add? &
Q1: shared transfer seen, not yet separated from prevalence. Q2: no inferential framework yet. Q3: DR's gain depends on the estimand. &
Not yet inferential &
\pageref{core:q123} & 8 \\
What did we decide, and what runs next? &
New runs: patients with $\geq$100 malignant cells (P02, P03, P04, P09 today); re-weight the Stage~2 labels; downsampling and Yu-Hsiang's data come next. &
Decisions and plans; nothing run yet &
\pageref{core:next} & 9 \\
\bottomrule
\end{tabularx}
}

\vspace{6pt}
{\small
\noindent\textbf{What changed since the team report, in three lines.}
(1) The imbalance specific to P02--P04 sits mainly in the Stage~2 labels, not in the Stage~1 factors; the patient with the most Stage~1 leverage is P01 (pages~\pageref{core:shape}--\pageref{core:d26}).
(2) A zero in the atlas is often an L1 threshold, not an absence of signal (page~\pageref{core:pcafa}).
(3) PC8 is named by its poles, and its protein correlates are mostly cell identity (pages~\pageref{core:pc8}--\pageref{core:protein}).\par}
   % p1: title + answers at a glance
% Core page: the cohort and the yardstick (half page), plus the key terms. Numbers: full report section 1;
% per-patient counts from study_config.json, as in Figure 1 of the full report.
\coreq{The cohort and the yardstick}{core:cohort}{what are we classifying, and how do we score it? (primer)}

\begin{coreanswer}
60,118 CD34-sorted MRD CITE-seq cells from 13 patients after transplant; 1,569 (2.61\%) are malignant. Seven patients are \emph{supported} (at least 10 malignant cells and both classes): P01, P02, P03, P04, P06, P07, P09. We report \emph{lift}: held-out AUPRC minus the patient's malignant prevalence $\pi$, in percentage points (pp), because a constant score has AP exactly $\pi$.
\end{coreanswer}

{\footnotesize
\setlength{\tabcolsep}{4pt}
\noindent\begin{tabularx}{\linewidth}{@{}l *{7}{>{\centering\arraybackslash}X} >{\centering\arraybackslash}p{2.9cm}@{}}
\toprule
 & P01 & P02 & P03 & P04 & P06 & P07 & P09 & other 6 (training only) \\
\midrule
Cells & 11,860 & 1,462 & 2,327 & 1,128 & 3,259 & 2,185 & 17,126 & 20,771 \\
Malignant cells & 45 & 211 & 758 & 102 & 44 & 22 & 375 & 12 (P05 8, P13 4; 4 with 0) \\
Prevalence $\pi$ & 0.38\% & 14.4\% & 32.6\% & 9.0\% & 1.4\% & 1.0\% & 2.2\% & -- \\
\bottomrule
\end{tabularx}\par}

\vspace{3pt}
\evidence
\begin{itemize}[leftmargin=1.3em]
\item P03, P02 and P04 are the \emph{highest-prevalence} patients (not the ones with the most malignant cells: P09 has more than P02 or P04). Together they hold 68.3\% of the malignant cells but only 8.2\% of all cells; P09 and P01 alone are 48.2\% of the cells.
\item Lift reads the gap above no skill (figure): P03 has the highest AP (46.0\%) but less lift ($+13.40$\pp) than P04 ($+16.27$) or P02 ($+15.94$), because its $\pi$ is 32.6\%.
\item An L1 probe that keeps no factor (\emph{all-zero}) scores AP $=\pi$ by construction: lift 0 then means ``nothing entered'', not ``no signal'' (P01 and P09 in the figure; page~\pageref{core:pcafa}).
\item Subtracting $\pi$ removes the floor, not the association: across 643 representations the median Spearman $\rho$ between lift and prevalence is 0.821.
\end{itemize}

\begin{figure}[H]
\centering
\includegraphics[width=0.66\linewidth]{figures/fig01b_yardstick.pdf}
\caption{Lift is the vertical gap between the held-out AP (blue) and the no-skill AP $=\pi$ (gray). IFN-$\alpha$ PCA $K=40$, L1 at the atlas $C=10^{-4}$, transductive Stage~1, leave-one-patient-out.}
\end{figure}

\doesnot{malignant biology apart from genotype (the normal cells are 99.9\% donor-derived, so malignant vs normal is almost recipient vs donor), or recurrence risk (lift ranks cells). P07, P06 and P01 rest on 22, 44 and 45 malignant cells.}

\begin{explainbox}{Key terms used on every page}
\footnotesize
\textbf{Stage~0}: a gene panel (e.g.\ hallmark IFN-$\alpha$ response, 92 genes; 41 biological panels). \textbf{Stage~1}: dimension reduction (DR) of the panel's genes into $K$ factors (PCA, factor analysis FA, FactoSig, FactoSig-promax; $K$ = 5, 10, 20, 40), fit without labels. \textbf{Stage~2}: a penalized logistic regression (\emph{probe}) on the factor scores; L1, L2 or elastic net (EN), penalty strength set by $C$ (smaller $C$, stronger penalty). \textbf{Representation}: one panel $\times$ method $\times$ $K$, e.g.\ IFN-$\alpha$ PCA40; the \emph{atlas} scores all of them. \textbf{LOPO}: leave one patient out; train on the others, score the held-out patient. \textbf{Transductive}: Stage~1 was fit on all cells, the held-out patient's included (labels never); \textbf{inductive}: fit without them. \textbf{AUPRC} (area under the precision--recall curve) is computed as average precision, \textbf{AP}. \textbf{Median / P20 lift}: the median and 20th percentile of lift over the seven supported patients; \textbf{pooled AP}: AP over all their held-out cells; three \textbf{estimands} (quantities estimated) that can rank methods differently. \textbf{Leverage test}: refit Stage~1 without a patient set and measure how far the malignant axis moves, against 20 random drops of as many cells and malignant cells; \textbf{matched $|r|$}: correlation of the original factor with its best match after the refit (1 = unchanged); \textbf{contrast sets}: P01, P06, P09 and P01+P09, the comparison sets that are not highest-prevalence. \textbf{CLR}: centred log-ratio normalized protein counts (no dsb background correction). \textbf{Descriptive}: an effect size on these data, with no calibrated test and no patient-level interval.
\end{explainbox}

\fullref{section 1, \emph{The cohort and the yardstick}; Appendix A (glossary).}
   % cohort and yardstick (+ key terms)
% Core page: do the highest-prevalence patients shape the factors? IFN-alpha part. Numbers: full report
% section 2 (D16, D17, D28). Figure: fig02a panels (a) and (c), cropped by trim (fig02a is 449 x 383 bp:
% panel a spans the top 145 bp, panel c the bottom 93 bp; recheck if figs_02_shape.py changes the layout).
\coreq{Do P02, P03 and P04 shape the factors? (IFN-\texorpdfstring{$\alpha$}{alpha})}{core:shape}{``patients with more malignant cells (P02, P03, P04) disproportionally shape the factors''; re-weight each patient equally in Stage~1? (team report)}

\begin{coreanswer}
\textbf{Not more than other patients do, on the panels tested.} Refitting the IFN-$\alpha$ Stage~1 basis without them leaves the malignant axis almost unchanged, and on all 41 panels they hold about their cell share of the factor variance. Where they dominate is Stage~2: 68.3\% of the malignant labels. MX chose to re-weight the Stage~2 labels (2026-10-10), in the runs with the new patients.
\end{coreanswer}

\evidence
\begin{itemize}[leftmargin=1.3em]
\item \emph{Two different imbalances} (panel a): P02+P03+P04 are 8.2\% of the cells, which is what the label-blind Stage~1 sees, but 68.3\% of the malignant cells, which is what Stage~2 learns from. The cell-heavy patients are P09 (28.5\%) and P01 (19.7\%).
\item \emph{Score energy} (how much of the factor-score variance sits in a patient's cells, divided by its share of cells; 1 = proportional): over all 320 stored PCA/FA bases the group's median ratio is 0.96 (range 0.58--1.14), 1.04 after adjusting for cell type. IFN-$\alpha$ PCA40 is near the top at 1.11 (panel c).
\item \emph{Leverage test} (refit Stage~1 without them and measure how far the malignant axis moves, against 20 random drops of as many cells \emph{and} as many malignant cells): on IFN-$\alpha$ under 1\% of the axis's variance is lost (matched $|r| \geq 0.988$, the correlation with the best refit factor), above all 20 drops in 3 of 6 method $\times$ $K$ cells. The rule needs 4; on this panel it is applied post hoc.
\item \emph{Patient-equal Stage~1 weights} (pilot): P02--P04 go from 8.2\% to 23.1\% of the fit weight; PC8 rotates (matched $|r|$ 0.90); median per-patient change $+1.35$ (PCA40) and $+0.07$\pp{} (FA40) under tuned L1, but FA40 loses a median 1.13\pp{} under L2.
\item \emph{P01 has the most leverage} of any set (IFN-$\alpha$ PCA40: 15$\times$ its controls). Its extra score energy is mostly its cell-type mix (ratio 1.24, 1.03 adjusted); its leverage is not explained.
\end{itemize}

\begin{figure}[H]
\centering
\includegraphics[width=0.80\linewidth,trim=0 238 0 0,clip]{figures/fig02a_who_dominates.pdf}\\[2pt]
\includegraphics[width=0.80\linewidth,trim=0 0 0 290,clip]{figures/fig02a_who_dominates.pdf}
\caption{(a) Each patient's share of all cells (light; its weight in the unweighted Stage~1 fit) and of the malignant cells (dark; its share of the Stage~2 positive labels). (c) The P02+P03+P04 score-energy ratio on each of the 320 stored PCA and FA bases (41 panels $\times$ $K$ 5/10/20/40); IFN-$\alpha$ PCA40 and FA40 circled. Panel (b) of the full report's figure (per patient) is left out here.}
\end{figure}

\doesnot{why a set moves an axis (a leave-out refit is an in-silico perturbation), FactoSig bases, or leverage on the other 38 panels; bases are transductive.}

\fullref{section 2, \emph{Do the highest-prevalence patients shape the factors?} (boxes ``Score energy'', ``The leverage test, step by step'', ``How patient-equal re-weighting works''). \src{2026-10-08\_d16-factor-patient-matrix; 2026-10-08\_d17-composition-matched-controls; 2026-10-09\_d28-celltype-score-shares}}
         % do P02/P03/P04 shape the factors? (IFN-alpha; D16, D17, D28)
% Core page: allograft rejection and apoptosis, the leverage test (D26). Numbers: full report section 2b.
% Figure: fig02b_d26a panels (a-c), cropped by trim (fig02b_d26a is 479 x 386 bp; panels a-c, their x label
% and the title span the top 220 bp, so trim removes the bottom 165 bp; recheck if figs_02b_d26.py changes the layout).
% Numbers not in the D26 done entry (PCA-only and FA-only ranges, control |r| ranges, normal-only refits, the
% 17-of-18 count) are recorded in devlog 2026-10-10_core-report-v3 (Results, with their source tables).
\coreq{Allograft rejection and apoptosis: what did we do?}{core:d26}{``What did we do with allograft and apoptosis? Please update all that as a new section in the report.'' (MX, 2026-10-10, A-20261010-13; also A-20261009-02, -15)}

\begin{coreanswer}
\textbf{We ran the IFN-$\alpha$ leverage test on both panels (D26, MX's screen run of 10 October). By the rule fixed before the run, P02, P03 and P04 do not shape these factors more than other patients do.} Leaving them out moves the malignant axis beyond random drops of the same composition, mostly through P03, but on the primary axis never beyond what leaving out another patient set does; P01 moves it most. So hypothesis H-S1-05 (``they shape the factors'') is not supported for PCA and FA.
\end{coreanswer}

\evidence
\begin{itemize}[leftmargin=1.3em]
\item \emph{Design} (as for IFN-$\alpha$, on the two programs where DR gains most): refit Stage~1 without P02+P03+P04, without each of them, or without a contrast set (P01, P06, P09, P01+P09); PCA and FA at $K = 10$, 20, 40; 2,364 refits. Rule: (1) \emph{beyond composition} if the axis moves more than all 20 composition-matched drops in at least 4 of 6 method $\times$ $K$ cells; (2) \emph{support} also needs the group to beat every contrast set in at least 4 of 6.
\item \emph{(1) is met} (panels b, c): above all 20 drops in 5 of 6 cells (allograft) and 6 of 6 (apoptosis), at 1.3--4.1$\times$ and 2.0--4.0$\times$ the drops' mean (IFN-$\alpha$, panel a: 3 of 6). P03 alone accounts for 74--93\% and 72--88\% of what the three together lose.
\item \emph{Size} (rule (4): matched $|r| \geq 0.98$ is ``small''): small for apoptosis ($|r| \geq 0.994$) and allograft FA ($\geq 0.988$). Allograft PCA's primary axis (its PC8, not the IFN-$\alpha$ PC8 of page~\pageref{core:pc8}) is not small: $|r|$ 0.945--0.970, 0.9--2.8\% of its variance lost. But the 20 random drops of 4,917 cells, 1,071 of them malignant, rotate it about as much ($|r|$ 0.90--0.97): that axis is partly made by the malignant cells, whoever they belong to.
\item \emph{(2) fails}: the group beats every contrast set in 0 of 6 primary-axis cells on each panel (1 of 6 on the apoptosis probe axis). P01 (11,860 cells, 45 malignant) moves the axes 21--141$\times$ and 48--128$\times$ its own controls, the largest ratio in 17 of 18 primary-axis cells over the three panels.
\item \emph{What it means}: it fits MX's choice to re-weight the Stage~2 labels. For healthy-cell downsampling, P01 (about 263 healthy cells per malignant cell, the most of the supported patients) is where thinning the Stage~1 fit would bite; a prediction, not a result.
\end{itemize}

\begin{figure}[H]
\centering
\includegraphics[width=0.86\linewidth,trim=0 165 0 0,clip]{figures/fig02b_d26a_ratio_by_set.pdf}
\caption{Each column is one patient set left out of the Stage~1 fit (02 = P02; 02+03+04 = the three together); each marker is one method $\times$ $K$ (blue PCA, orange FA; triangle $K$ = 10, square 20, circle 40; open = inside its 20 drops). The y-axis is how far the primary axis moves (axis loss) divided by the mean of that set's 20 composition-matched drops; above 1 = more than random drops. The number over each column counts cells above all 20; P01's column is shaded. The full report's panels (d--f) show matched $|r|$.}
\end{figure}

\doesnot{FactoSig or FactoSig-promax (the atlas winners on both panels) or the other 38 panels; held-out lift (no Stage~2 probe was refit); why P01 moves the axes (proposed: a cell-type-matched control, W2.9).}

\fullref{section 2b, \emph{What did we do with allograft rejection and apoptosis?} (rule table, next steps). \src{2026-10-09\_d26-leverage-allograft-apoptosis; 2026-10-10\_d26-leverage-allograft-apoptosis-done; 2026-10-10\_core-report-v3}}
 % allograft rejection and apoptosis: the leverage test (D26)
% Core page: should we aggregate patients? (D14). Numbers: full report section 3.
\coreq{Should we aggregate patients?}{core:aggregate}{should we aggregate patients, given the per-patient factor distributions? (team report)}

\begin{coreanswer}
\textbf{Yes, on the current Stage~1 bases.} The patients' cells land in the same factor-score ranges: of 47 representations (41 biological panels and 6 broad gene-bag controls), 45 pass all four overlap gates, 2 are \emph{location shifts} (one patient's cells sit shifted) and none looks like a patient-specific subtype.
\end{coreanswer}

\evidence
\begin{itemize}[leftmargin=1.3em]
\item \emph{Overlap}: the shared area of two patients' histograms of one factor's scores (1 = identical, 0 = disjoint), averaged over factors with variance weights, for normal and malignant cells separately.
\item \emph{The gates} (a fixed definition of ``aggregate''; judgment calls, not calibrated): the median patient pair shares at least 0.75 of its normal-cell and 0.45 of its malignant-cell histogram; no patient's normal cells share less than 0.60 with the rest; the malignant overlap is less than 0.15 below that of label-shuffled patients.
\item On the 41 biological panels the median pair shares 0.777--0.890 of its normal-cell histogram (median 0.843) and 0.538--0.657 of its malignant-cell histogram (median 0.610).
\item The two location shifts: allograft rejection FactoSig-promax $K = 40$ (P02's normal cells overlap the rest at 0.575) and the HVG 10k control bag. The only patient-private axis, P03's, appears in three broad gene bags and in none of the 41 biological representations.
\end{itemize}

\begin{figure}[H]
\centering
\includegraphics[width=0.98\linewidth]{figures/fig03a_decision_map.pdf}
\caption{The aggregation call for every representation: one point per Stage~0 panel at its atlas-selected method and $K$ (colour = method; open squares = broad gene-bag controls; gray dots = label-shuffled patients). The shaded region passes the overlap gates; the red band would flag a subtype. The title counts the 41 biological panels (40 pass); with the 6 controls it is 45 of 47.}
\end{figure}

\doesnot{how much each patient \emph{contributes} to a factor (that is the leverage test) or whether the same factors drive lift; the basis is transductive and the gates are judgment calls (permutation $p$ at its floor). P07, P06 and P01's malignant overlaps rest on 22--45 cells.}

\fullref{section 3, \emph{Should we aggregate patients?} (worked example on PC8, the three calls drawn as shapes). \src{2026-10-02\_d14-patient-factor-overlap}}
          % should we aggregate patients? (D14)
% Core page: IFN-alpha PCA40 vs FA40, the L1 entry rule (D19). Numbers: full report section 4.
\coreq{IFN-\texorpdfstring{$\alpha$}{alpha}: why does PCA40 reach four patients but FA40 only two?}{core:pcafa}{why does PCA40 transfer to four patients but FA40 only to P03 and P04? Hypothesis (a): FA keeps allocating factors to the highest-prevalence patient. (team report)}

\begin{coreanswer}
\textbf{A Stage~2 regularization threshold, not missing FA structure.} Under L1 a factor can enter a fold's probe only if its standardized class-mean gap exceeds $4/(C\,n_{\text{train}})$ ($n_{\text{train}}$ = training cells). FA40's best factor, FA7, misses that bar by 1.4--3.3\% in P02, P06 and P07, so those folds keep no factor and score lift 0 by construction. With $C$ tuned per fold, FA40 has positive lift in all 7 supported patients. Hypothesis (a) is retired. (The team's ``four'' are PCA40's patients above 5\pp; its fifth positive patient, P06, is at $+0.06$\pp, which the figure counts as lift $>$ 0.)
\end{coreanswer}

\evidence
\begin{itemize}[leftmargin=1.3em]
\item \emph{The bar moves with the fold}: at the atlas setting (L1, $C = 10^{-4}$) PC8 clears it by 7.3--8.5\% in P02, P06 and P07, FA7 misses it by 1.4--3.3\%. A 5\% higher $C$ already gives FA40 18.06\pp{} in P02 and 6.97\pp{} in P07 (figure, second panel).
\item \emph{Tuned $C$}: positive lift in 7 of 7 patients for both bases under L1, L2 and elastic net; tuned-L1 median lift 8.71\pp{} (PCA40) and 7.03\pp{} (FA40).
\item \emph{The rule is exact}: it reproduces 14/14 IFN-$\alpha$ folds and 2,772/2,772 audited atlas folds. Across the atlas, 528 of 783 cells select the lowest $C$ on the grid and 136 folds in 49 cells sit at the knife-edge of entry, so a zero in the atlas is often a threshold, not an absence of signal.
\item \emph{Hypothesis (a)}: FA40 does not favour the highest-prevalence patients more than PCA40 does. P02--P04 hold 8.4\% of FA40's and 9.1\% of PCA40's score energy against 8.2\% of the cells, and P03 tops none of either basis's factors.
\end{itemize}

\begin{figure}[H]
\centering
\includegraphics[width=0.93\linewidth]{figures/fig04b_policy_lift.pdf}
\caption{Held-out lift per supported patient for PCA40 (blue) and FA40 (orange) under four $C$ policies; hollow = all-zero probe (lift 0 by construction). Look at the hollow orange circles for P02, P06 and P07 in the first panel, and how they fill in once $C$ rises by 5\% or is tuned.}
\end{figure}

\doesnot{that PCA is better or worse than FA (nothing here ranks them); one panel at $K = 40$; the atlas $C$ was chosen with the held-out patients' outcomes (not nested).}

\fullref{section 4, \emph{IFN-$\alpha$: why does FA40 reach only P03 and P04?} (entry-margin figure; box on the entry rule). \src{2026-10-08\_d19-regularization-contract}}
          % IFN-alpha PCA40 vs FA40: the L1 entry rule (D19)
% Core pages: what PC8 is (D18), then the surface-protein page (D27). Numbers: full report section 5, and,
% for MX's 2026-10-10 protein questions (A-20261010-01, -03..-06), devlog 2026-10-10_protein-association-for-pavan-v2.
% Both figures are placed at their drawn width (6.5 in, figs_05_pc8.py), so their text prints at its drawn size.
\coreq{What is PC8?}{core:pc8}{PC8's opposite-sign gene groups: what is it? (team report)}

\begin{coreanswer}
\textbf{One axis that contrasts two groups of IFN-$\alpha$ genes; name it by its poles.} Oriented by its coefficient in the L1 probe ($C=10^{-4}$), malignant cells score toward TXNIP, IFITM2, TENT5A (the team's ``negative'' group) and normal cells toward IFI27, PSME2, PSMA3. It carries held-out lift in five of seven patients, but it is not identifiable as ``factor 8'': a refit of Stage~1 can spread the same direction over several factors.
\end{coreanswer}

\evidence
\begin{itemize}[leftmargin=1.3em]
\item \emph{Lift}: PC8 alone gives 5.47--16.27\pp{} of held-out lift in P02, P03, P04, P07 and P09 (FA7 alone 4.81--19.87\pp), and under 1\pp{} in P01 and P06.
\item \emph{Poles} (figure): the malignant pole runs from TXNIP to ELF1, the normal pole from IFI27 to USP18. Interferon-stimulated genes sit on both poles, so the axis contrasts two parts of the response rather than measuring ``more interferon''; the team's two gene lists are its two ends.
\item \emph{FA7 is related, not identical}: its loadings resemble PC8's (similarity 0.69 on a 0--1 scale, below the 0.85 at which two factors are usually called fairly similar), and it has the same sign as PC8 on 19 of the 20 pole genes in the figure.
\item \emph{Not identifiable} (a refit can return a different single direction): PC8's variance is within 5.5\% of PC7's; without P01 in Stage~1, PC8's best single match falls to $|r|$ 0.58 while 98.8\% of it stays in the refit basis; FA7 is only loosely pinned down by the data (determinacy 0.649, below the usual 0.707 line).
\end{itemize}

\begin{figure}[H]
\centering
\includegraphics[width=6.5in]{figures/fig05a_pc8_poles.pdf}
\caption{PC8's poles, oriented toward malignant: the top 10 genes per pole (bars = PC8 weight; dots = FA7 loading; bold = named in the team report).}
\end{figure}

\doesnot{the pole genes as differential expression (they are loadings on one axis); ``factor 8'' as a stable object across refits; malignancy apart from recipient origin or batch.}

\fullref{section 5, \emph{What is PC8} (identifiability box and figure; the 1,810 factor cards). \src{D18; 2026-10-08\_d18-factor-cards}}

\coreq{Have we integrated the protein data? Does it track PC8?}{core:protein}{Pavan: ``have you been able to integrate the protein data?''; MX: how are isotypes known, and why do they define noise? Is the CD34 correlate the CD34 sort? What other surface protein tracks with healthy? What does ``19 out of 20'' mean? (A-20261010-01, -03..-06)}

\begin{coreanswer}
\textbf{Not into the factors: Stage~1 is fit on RNA only. We use the 170 surface proteins to read the factors.} PC8's protein correlates are mostly cell identity: CD34 tracks PC8 across cell types, not within them. Inside one patient and one cell type, malignant cells still score higher on PC8 than healthy cells in 19 of 21 comparisons.
\end{coreanswer}

\evidence
\begin{itemize}[leftmargin=1.3em]
\item \emph{CD34} (normal cells, Spearman $\rho$ with oriented PC8): median $+0.245$ within patients (7/7 same sign) but $+0.019$ inside predicted cell types (panel a). Every cell passed a CD34+ sort, yet CD34 still varies with differentiation stage (brightest in HSC/MPP), so the correlate says which cell type a cell is.
\item \emph{Other healthy-cell proteins} shrink the same way inside cell types: CD131 $+0.190$, CD112 $+0.154$, CD81 $-0.201$, CD22 $-0.246$ within patients (6--7 of 7).
\item \emph{Malignant cells} (rule for a \emph{robust} correlate: median within-patient $|\rho| \geq 0.20$ and $\geq$5 of 7 patients with the same sign): CD99 (pooled $+0.222$) is $+0.159$ within patients and fails it (panel b); CD48 (pooled $-0.192$) is the robust one, $-0.308$ within patients, 6 of 7 (panel c), because P03, with half the malignant cells, sits near 0 and dilutes the pooled value.
\item \emph{``19 of 21''}: for each supported patient $\times$ predicted cell type with at least 10 cells of each label, the malignant-minus-healthy mean PC8. Malignant is higher in 19 of the 21 pairs (the first Pavan figure counted 19 of 20 pairs, over fewer cell types). PC8's sign was set by the malignant-vs-healthy probe, so only the persistence inside one patient and cell type is new.
\item \emph{Isotypes}: 9 of the 170 protein columns are named \texttt{Isotype-<clone>} in the data: control antibodies that bind nothing specific. Their correlation with a factor shows how large a $\rho$ arises from background staining alone, so we use it as a noise floor (gray band); a design choice, not a test.
\end{itemize}

\begin{figure}[H]
\centering
\includegraphics[width=6.5in]{figures/fig05e_rho_anchors.pdf}
\caption{One pooled $\rho$ taken apart: per patient (dots; cell counts in parentheses), the median of the seven, pooled over all cells, and inside predicted cell types; gray band = isotype range; dotted lines = the robust rule's $|\rho| = 0.20$.}
\end{figure}

\doesnot{any protein as a malignancy marker; dsb-normalized protein (CLR only); a factor fit that uses the protein (not attempted); malignancy apart from recipient origin or batch.}

\fullref{section 5, \emph{What is PC8} (D27 box ``How to read $\rho$'', every factor card); Pavan's v2 figures and captions: \texttt{research/\allowbreak team\_updates/\allowbreak 2026-10-09\_protein\_association\_for\_pavan/}. \src{D27; 2026-10-09\_d27-factor-modality-assoc; 2026-10-10\_protein-association-for-pavan-v2}}
        % what PC8 is, and surface protein (D18, D27)
% Core page: Option 2, excluding the held-out patient from Stage 1 (D17). Numbers: full report section 6.
\coreq{Excluding the held-out patient from Stage~1}{core:inductive}{Option~2: fit Stage~1 without the held-out patient, so the factors never see it (team report)}

\begin{coreanswer}
\textbf{Yes, on the one panel tested (IFN-$\alpha$; PCA and FA at $K$ = 10, 20, 40).} With Stage~1 fit without the held-out patient and that patient's cells projected onto the frozen factors (\emph{inductive}), held-out lift stays positive in all 7 supported patients under tuned $C$ or L2 (126 of 126 patient $\times$ setting pairs). The transductive numbers are slightly optimistic per patient, by little.
\end{coreanswer}

\evidence
\begin{itemize}[leftmargin=1.3em]
\item \emph{Cost per patient}: against the transductive refit, the median per-patient change is between $-1.29$ and $-0.05$\pp{} across the 18 tuned/L2 settings; always negative, but small next to median lifts of 2.5--10.6\pp.
\item \emph{Pooled AP is more sensitive}: FA40 under tuned L2 falls from 0.1325 to 0.0794.
\item \emph{The axis survives}: the discriminative factors (PC8, FA7) survive the refit, and a basis fit on normal cells only carries them too (median change $+0.02$\pp{} for PCA40 and $+0.25$\pp{} for FA40 under tuned L1; figure).
\item \emph{Single patients move more}: FA40 P03 $-3.83$\pp, PCA10 P04 $+6.13$\pp; with seven patients these are not identifiable effects.
\end{itemize}

\begin{figure}[H]
\centering
\includegraphics[width=\linewidth]{figures/fig06a_inductive_lift.pdf}
\caption{Held-out lift per supported patient at $K = 40$, $C$ tuned per fold, L1, on three Stage~1 bases: transductive refit (open circle), inductive (filled circle) and normal cells only (diamond).}
\end{figure}

\doesnot{the other representations (the 47 atlas-selected ones, FactoSig included, are untested); fully inductive tuning (the inner folds reuse the outer fold's basis); new data such as Yu-Hsiang's patients.}

\fullref{section 6, \emph{Option 2: inductive Stage 1} (all 18 settings, paired changes). \src{2026-10-08\_d17-basis-variant-pilot}}
          % excluding the held-out patient from Stage 1 (D17)
% Core page: learning curves (D15, D15b). Numbers: full report section 7 and its figures.
\coreq{Learning curves: AUPRC against cells and patients}{core:curves}{plot AUPRC against the number of cells or patients; is it useful for the grant? (team report)}

\begin{coreanswer}
\textbf{We can draw them, but they do not yet justify sequencing more patients.} With $C$ retuned for every training set, one supported training patient already gives positive lift (median 4.45--8.79\pp; the six unsupported patients, with fewer than 10 malignant cells (P05, P08, P10--P13), are always in training). Going from two to six supported training patients adds a median 0.30--2.34\pp{} per held-out patient, with a 95\% interval that includes 0 for three of four representations.
\end{coreanswer}

\evidence
\begin{itemize}[leftmargin=1.3em]
\item \emph{Why $C$ is retuned}: at fixed $C$ the L1 entry bar $4/(C\,n_{\text{train}})$ rises as cells are removed, so a small training set is a stronger penalty; at the atlas $C$ the IFN-$\alpha$ curve is 0 up to 25,000 training cells.
\item \emph{Cells} (same 12 training patients, $C$ fixed): median apoptosis lift rises from 7.65\pp{} at 500 training cells to 11.52\pp{} on the full fold.
\item \emph{Patients} (all 63 training subsets per held-out patient; 1,764 units over the four representations): median gain from two to six patients, with 95\% patient-bootstrap interval: IFN-$\alpha$ PCA40 $+1.67$ [$-0.44$, 2.89], FA40 $+0.30$ [$-0.53$, 4.18], apoptosis $+1.14$ [$-0.53$, 2.38], allograft $+2.34$ [0.14, 2.73]. The PCA40 value replaces the 2.16 [0.20, 2.48] of the 2026-10-08 team update.
\item \emph{For the grant}: as feasibility, yes (transfer works with one supported training patient); as evidence that more patients raise AUPRC, not yet: seven patients are too few to measure scaling.
\end{itemize}

\begin{figure}[H]
\centering
\includegraphics[width=0.95\linewidth]{figures/fig07b_patient_count.pdf}
\caption{Patient-count curve: median held-out lift over the seven held-out patients by the number $k$ of supported training patients, $C$ and penalty retuned per training set (coloured), with its 95\% patient-bootstrap interval (band); gray = fixed atlas $C$.}
\end{figure}

\doesnot{that sequencing more patients raises AUPRC (nothing is extrapolated); an inductive curve (the basis saw every patient); patients apart from their malignant cells (adding a patient adds its malignant cells).}

\fullref{section 7, \emph{Learning curves} (cell-depth curve, per-patient gains). \src{2026-10-02\_d15-cell-count-learning-curve; 2026-10-08\_d15b-full-design}}
    % learning curves (D15, D15b)
% Core page: the three scientific questions (Q1-Q3) and the claim boundary. Numbers: full report section 8.
\coreq{The three scientific questions, and what we can claim}{core:q123}{Q1 is the signal shared or patient-specific? Q2 how strong is the evidence? Q3 what does DR add over direct genes? (team report)}

\begin{coreanswer}
\textbf{Q1.} Shared transfer is seen, but not yet separated from prevalence. \textbf{Q2.} Effect sizes are well defined; there is no inferential framework yet. \textbf{Q3.} It depends on the estimand and on how the winning DR representation is chosen; large gains are confined to a few programs. Everything is descriptive, on a transductive Stage~1 basis, over seven supported patients.
\end{coreanswer}

\evidence
\begin{itemize}[leftmargin=1.3em]
\item \emph{Q1, shared or patient-specific}: tested directly on IFN-$\alpha$ only, where one axis (PC8, or the related FA7) carries held-out lift in five of seven patients, including low-prevalence P07 and P09. Across the 640 biological DR atlas cells, lift above 0 reaches every supported patient, but lift above 5\pp{} occurs in 64--80\% of cells for P02, P03, P04 and in at most 23\% for the other four (figure).
\item \emph{Q2, strength of evidence}: lift above the constant-score no-skill AP $=\pi$ is well defined, but there is no calibrated null, no patient-level interval on the atlas and no nested selection (the best representation is picked on the folds it is scored on).
\item \emph{Q3, DR vs direct genes}: when the best DR cell is picked by pooled AP, it beats a probe on the panel's genes directly in 38/41 programs on pooled AP, 16/41 on median lift and 17/41 on P20 lift; picked by median lift, it wins on median lift in 28/41. Large gains: allograft rejection, apoptosis, IFN-$\alpha$; broad control bags favour direct genes.
\item \emph{Claim boundary}: always quote pooled AP, median lift and P20 lift together. The study's claim list is THEME \S6; no result here is inferential or ready for publication.
\end{itemize}

\begin{figure}[H]
\centering
\includegraphics[width=0.92\linewidth]{figures/fig08a_patient_concentration.pdf}
\caption{Where the lift sits: for each supported patient (ordered by prevalence $\pi$), the share of the 640 biological DR atlas cells with held-out lift above 0 (light) and above 5\pp{} (dark).}
\end{figure}

\doesnot{inductive deployment (Stage~1 saw every patient; LOPO is label transfer), causal biology, recurrence risk (lift is a cell-level ranking), or significance; best-DR win counts will shrink under nested selection.}

\fullref{section 8, \emph{The three scientific questions} (variance estimands, DR-vs-direct table and figure). \src{D1, D1e, D1f, D6, D9, D12--D19; 2026-09-11\_d1-atlas-783-d1e-d1f; 2026-08-31\_d6-d9-meeting-followup}}
    % Q1-Q3 and the claim boundary
% Core page: decisions taken (2026-10-09/10) and what runs next with the new patient data. Numbers: full
% report section 9 (Table "What each inclusion threshold implies") and sections 2/2b; decisions recorded in
% devlog 2026-10-10_decisions-new-run-inclusion-reweighting-lfs (asks A-20261010-07..10).
\coreq{Decisions taken, and what runs next with the new data}{core:next}{``Let's do the $>$=100 malignant cells. Reweight should target stage 2 label \ldots{} These all sounds like decisions relevant for runs when we have the extra patient data.'' (MX, 2026-10-10)}

\begin{coreanswer}
All four items below apply to the runs on Yu-Hsiang's new patients; existing results keep the seven supported patients. Nothing new has been computed for them yet.
\end{coreanswer}

{\small
\renewcommand{\arraystretch}{1.2}
\noindent\begin{tabularx}{\linewidth}{@{}>{\raggedright\arraybackslash}p{3.2cm} >{\raggedright\arraybackslash}X@{}}
\toprule
\textbf{Decision or plan} & \textbf{What it means} \\
\midrule
\textbf{Only patients with $\geq$100 malignant cells} (decided 2026-10-10; W8.5) &
On the current cohort: P02, P03, P04 and P09 (any cut from 46 to 102 gives the same four). Patients below 100 are left out of every step (Stage~1 fit, Stage~2 training, summaries); that is our reading of MX's note, to confirm when the plan is locked (W8.2). Patients without malignant cells (P08, P10, P11, P12) are out too. \emph{Consequence}: three of the four are the highest-prevalence patients, so the same probe reads better without being better: median lift over the four is 1.8--2.7$\times$ today's seven-patient median (IFN-$\alpha$ PCA40 15.52 vs 8.71\pp; P20 lift 11.54 vs 1.37\pp). $C$ is tuned on 3 inner folds instead of 6, and the smallest exact sign-test $p$ is 1/16, so it cannot reach 0.05. Summaries must name their patient set. \\
\addlinespace
\textbf{Re-weighting targets the Stage~2 labels} (decided 2026-10-10; W2.3) &
Each training patient's malignant cells get the same total weight in the logistic loss, and likewise its normal cells; $C$ is retuned per condition. It runs as the Stage-2-only sample-weight scheme of the downsampling design. No re-weighted Stage~1 refit is planned (pages~\pageref{core:shape}--\pageref{core:d26}). \\
\addlinespace
\textbf{Healthy-cell downsampling} (plan drafted, not locked; H-S2-09) &
Tests whether the healthy-to-malignant \emph{ratio} (about 2 healthy cells per malignant cell in P03, about 263 in P01), not only the malignant count, makes classification hard. Schemes compared at equal kept cells (random, proportional within patient $\times$ cell type $\times$ pseudotime, type-equalizing, patient-capped); malignant retention $\{1, \frac12, \frac14\}$ $\times$ healthy $\{1, \frac12, \frac14, \frac18\}$; $C$ retuned; success and failure fixed in advance. The earlier donor downsampling (D10, Stage~2 only, $C$ fixed) gave no gain under L2. Cost about 100--150 CPU-hours at seven held-out patients (it scales with their number: 4 plus the new patients under the 100-cell rule); the runner (W8.6) is not written. \\
\addlinespace
\textbf{Yu-Hsiang's new patients} (W8.1) &
The only external test: patients that took no part in choosing panel, basis, probe or $C$. The email and a pre-registered integration plan are drafted; the email is not sent. FactoSig bases cannot yet project new cells, so each FactoSig-promax choice has a PCA/FA fallback. \\
\bottomrule
\end{tabularx}\par}

\vspace{4pt}
\noindent\textbf{Open decisions for the study lead (MX)}: send the Yu-Hsiang email (W8.1); at the plan lock, confirm that patients below 100 malignant cells are left out rather than kept in training (W8.2); a go-ahead for the leverage test on all biological panels (W2.8, about 3 hours in a screen session).
\textbf{Repository}: the Git LFS lines are removed (2026-10-10, W0.2). This core PDF ships with the repository (a \texttt{.gitignore} exception; it reaches a clone with the push, W0.3); the full and team PDFs and all figures are gitignored and rebuilt with \texttt{build.sh}, which needs the experiment data (\texttt{README.md} in this directory; W0.15).

\doesnot{any result on the new patients, or that downsampling helps; the plans are drafts, and their channels and predictions are arguments, not results.}

\fullref{section 9, \emph{Remaining TODOs and decisions} (inclusion table, downsampling design, timing); plans \texttt{integration\_plan.md} and \texttt{healthy\_downsampling\_plan.md} in \texttt{workstreams/\allowbreak W8\_new\_cohort/}. \src{2026-10-10\_decisions-new-run-inclusion-reweighting-lfs; 2026-10-09\_w8-healthy-downsampling-plan}}
     % decisions taken and what runs next with the new data
% Not built (2026-10-10, version 3.1): a three-page appendix with the full report's surface-protein explainers
% (rho schematic, Table 3 + two checks, Figure 19). It would take the core to 15 pages, over the 12-page limit
% (A-20261010-11); pages 7-8 carry compact versions instead (Figure 6 bottom = fig05m; Table 1; Figure 7 = fig05l).
% To print the appendix anyway, uncomment the next line and restore the last row of the glance table on page 1.
% % NOT BUILT (2026-10-10, version 3.1): report_core.tex has this file's \input line commented out, because the
% appendix would take the core to 15 pages (limit 12). Pages 7-8 (06_pc8_protein.tex) carry compact versions
% instead: Figure 6 bottom (fig05m, Figure 19), Table 1 (Table 3's PC8 malignant rows) and Figure 7 (fig05l: the
% sign of rho and the two checks). Kept so the appendix can be printed on request (asked of MX, 2026-10-10).
% Core appendix (three pages after the 12-page body): the surface-protein explainer figures MX asked to see in
% the report (asks A-20261010-14..24; devlog 2026-10-10_protein-explainer-v3). The label core:appendix marks the
% first appendix page; build.sh counts the body pages as the pages before it (limit 12). Numbers: full report
% section 5 (Figures 19, 20, 22 and Table 3) and that devlog entry. Figures come from figs_05p_protein_explainer.py
% (which runs the Pavan update's make_figures_v3.py) and are placed at 6.5-6.8 in (drawn 7.0 in wide).
\coreq{Appendix A1. How is one $\rho$ computed, and what does its sign mean?}{core:appendix}{``PC8's cell score for all malignant cells is a vector of length (1569) \ldots{} spearman rho is just the correlation of the ranking of that with CD34's cell score or another surface protein's cell score.'' ``So the sign means how forward and reversed the rank of the factor cell score against the ADT protein cell score? We should include that in the schematics'' (MX, 2026-10-10; A-20261010-16, -23)}

\begin{figure}[H]
\centering
\includegraphics[width=6.5in]{figures/fig05k_rho_schematic_core.pdf}
\caption{From a factor and a protein to one number (panels A and D of full report Figure~20). (A) One PC8 score per malignant cell is a vector of 1,569 numbers; the CD99 levels (CLR) of the same cells, in the same order, are a second one. Spearman $\rho$ is the Pearson correlation of their ranks: $+0.222$. (B) Yes, the sign says whether the two rankings run forward (CD99: cells ranked high on oriented PC8 tend to rank high on CD99) or in reverse (CD48, $-0.192$); $|\rho|$ says how consistently. The full report's Figure~20 adds an 8-cell worked example and the same $\rho$ within patients (median $+0.159$, 6 of 7) and within patient $\times$ cell type (median $+0.150$, 9 of 11 groups).}
\end{figure}

\fullref{section 5, \emph{What is PC8}: Figure~20 and the box ``How to read $\rho$''. \src{2026-10-10\_protein-explainer-v3; D27}}

\coreq{Appendix A2. Table 3 in plain words, and two checks per protein}{core:appA2}{``The within (k/7), k means the number of patients that has the same rho sign, right? In type is a median value \ldots{} CD99 \ldots{} stayed as high as 0.150 \ldots{} But CD48 is somehow more robust'' and ``we should also include CD116$-$0.287, CD33$-$0.272, CD146 $+$0.270'' (MX, 2026-10-10; A-20261010-17, -18, -21)}

\begin{coreanswer}
\textbf{The robust rule reads only the within-patient column.} CD48, CD116, CD33 and CD146 pass it but fall to $|\rho| \leq 0.037$ inside one patient's cell type, so they mostly say which cell types sit at which end of PC8. CD99 fails it ($+0.159$, inside the isotype range) yet keeps $+0.150$ inside cell types (9 of 11 groups): a lead to check, not a finding. Yes, $k$ is the number of the 7 patients whose $\rho$ has the sign of the median.
\end{coreanswer}

\begin{table}[H]
\centering
\caption{PC8 against surface protein at three levels (the PC8 rows of full report Table~3, whose FA7 rows are in section~5). Spearman $\rho$ among cells of one label. In brackets: how many of the 7 patients, or of the patient $\times$ cell-type groups with $\geq 30$ cells (71 healthy; 11 malignant, from P02, P03, P04, P09), share the median's sign. Isotype reference: the largest $|\rho|$ of the 9 isotype controls in that column. $\ddagger$ added 2026-10-10.}
\footnotesize
\setlength{\tabcolsep}{6pt}
\renewcommand{\arraystretch}{1.02}
\begin{tabular}{@{}llrrr@{}}
\toprule
 & Protein & all cells pooled & median within patient & median within patient and cell type \\
\midrule
\multicolumn{5}{@{}l}{\emph{Healthy cells (58,549)}} \\
 & CD34 & $+0.251$ & $+0.245$ (7/7) & $+0.019$ (40/71) \\
 & CD22 & $-0.221$ & $-0.246$ (6/7) & $-0.046$ (52/71) \\
 & \emph{isotype reference} & \emph{0.056} & \emph{0.109} & \emph{0.043} \\
\multicolumn{5}{@{}l}{\emph{Malignant cells (1,569)}} \\
 & CD99 & $+0.222$ & $+0.159$ (6/7) & $+0.150$ (9/11) \\
 & CD48 & $-0.192$ & $-0.308$ (6/7) & $-0.011$ (6/11) \\
 & CD116 $\ddagger$ & $-0.108$ & $-0.287$ (7/7) & $-0.037$ (9/11) \\
 & CD33 $\ddagger$ & $-0.137$ & $-0.272$ (5/7) & $+0.026$ (6/11) \\
 & CD146 $\ddagger$ & $+0.090$ & $+0.270$ (6/7) & $+0.017$ (6/11) \\
 & CD63 $\ddagger$ & $-0.183$ & $-0.252$ (6/7) & $-0.138$ (9/11) \\
 & CD82 $\ddagger$ & $+0.138$ & $+0.238$ (7/7) & $+0.143$ (9/11) \\
 & GPR56 $\ddagger$ & $+0.203$ & $+0.137$ (6/7) & $-0.057$ (6/11) \\
 & \emph{isotype reference} & \emph{0.101} & \emph{0.227} & \emph{0.120} \\
\bottomrule
\end{tabular}
\end{table}

\begin{figure}[H]
\centering
\includegraphics[width=6.5in]{figures/fig05j_two_checks.pdf}
\caption{Two checks per protein (full report Figure~22): x = the median within-patient $\rho$, y = the median inside patient $\times$ cell type, one dot per protein (orange: isotypes). Gray bands: the isotype reach at each level; dotted: the robust rule's 0.20 size bar; dashed: $y = x$. Only CD63 and CD82 sit outside both bands in malignant cells; in healthy cells every protein has $|y| \leq 0.10$.}
\end{figure}

\coreq{Appendix A3. The cell-type counts and ``mean oriented PC8 score'' of full report Figure~19; FA7's two protein lists}{core:appA3}{``Why 643 HSC/MPP2? What's mean oriented PC8 score.'' ``So does that mean CD22 is less associated with normal cells?'' ``FA7 malignant cells' surface protein correlation is finding two different surface protein compare to its healthy counter part.'' (MX, 2026-10-10; A-20261010-14, -15, -19, -20)}

\begin{figure}[H]
\centering
\includegraphics[width=6.8in]{figures/fig05g_celltype_counts.pdf}
\caption{Full report Figure~19, relabelled: mean oriented PC8 score of healthy (open) and malignant (filled) cells per predicted type, all 13 patients pooled; (malignant / healthy) cell counts per type; right, the largest share of each type's malignant cells from one patient. Gray labels: malignant not higher.}
\end{figure}

\evidence
\begin{itemize}[leftmargin=1.3em]
\item \emph{The counts} are all cells of all 13 patients. Only the 10 types with at least 20 malignant cells are drawn: they hold 1,544 of the 1,569 malignant cells (pDC2 16, ncMono 6 and CLP 3 are left out), which is why the labels do not add up. HSC/MPP2 has 643 malignant cells (P03 438, P09 129, P02 45, P07 10, P01 9, P04 7, P06 4, P13 1) and 5,984 healthy ones; P03 alone supplies 68\% of the malignant ones and 92\% of the malignant NP1 cells, so a pooled gap can partly be a patient difference. Inside each patient and type, malignant is higher in 19 of 21 comparisons.
\item \emph{Mean oriented PC8 score}: each cell's PC8 score, z-scored over all 60,118 cells and multiplied by $-1$ so that positive points to the malignant-leaning end, averaged over the cells of one type and label; $+1$ = one cohort standard deviation.
\item \emph{CD22's negative sign} in healthy cells (PC8 $-0.246$, FA7 $-0.312$ within patients) does not mean CD22 is ``less associated with normal cells'': no malignant cell enters that $\rho$. It says healthy cells with more CD22 tend to sit toward the healthy-leaning end, mostly because CD22 is highest on predicted CLP and cDC, which sit there; inside cell types it shrinks to $-0.046$ and $-0.085$.
\item \emph{FA7's two lists} rank different cells: healthy cells are mostly CLP (21\%) and MEP-EP (18\%), malignant cells 41\% HSC/MPP2 and 18\% NP1. Each list holds the proteins that separate the cell types present, and they cross over weakly (only CD63 tracks FA7 in both labels).
\end{itemize}

\doesnot{any protein as a malignancy marker; a malignant-vs-healthy contrast in any $\rho$; dsb-normalized protein (CLR only); malignancy apart from recipient origin or batch.}

\fullref{section 5, \emph{What is PC8}: Figure~19, Table~3, paragraph ``Why FA7's normal-cell and malignant-cell proteins differ''. \src{2026-10-10\_protein-explainer-v3; D27; D18}}
 % appendix A1-A3 (D27)

\end{document}
